%% kapa: exam=1 round=27 year=2016 period=2 form=A subject=appraisal_law subject_ko=감정평가관계법규 seq=14 number=14
%% source: Q-Net 감정평가사 기출문제 (cst00309) artlSeq=5169402 file=extracted/1008/5169402/제27회 감정평가사 1차 문제/제27회 감정평가사 2교시문제(A형).pdf
%% conversion: kordoc 4.24.51 -> tools/gen1.py (auto); review=proofread-sonnet; flags=none
\begin{question}{14}
부동산가격공시 및 감정평가에 관한 법령상 과징금에 관한 설명이다. 다음 ( )안에 알맞은 것은?
\begin{framedtext}
국토교통부장관은 감정평가업자의 위반행위에 대해 업무정지처분을 하여야 하는 경우로서 그 업무정지처분이 표준지공시지가의 조사ㆍ평가 등의 업무의 정상적인 수행에 지장을 초래하는 등 공익을 해칠 우려가 있는 경우에는 업무정지처분에 갈음하여 감정평가사의 경우에는 ( ㄱ ) 이하, 감정평가법인인 경우는 ( ㄴ ) 이하의 과징금을 부과할 수 있다. \par
\end{framedtext}
\begin{choices}
  \item ㄱ: 5천만원, ㄴ: 5억원
  \item ㄱ: 5천만원, ㄴ: 10억원
  \item ㄱ: 7천만원, ㄴ: 7억원
  \item ㄱ: 1억원, ㄴ: 5억원
  \item ㄱ: 1억원, ㄴ: 10억원
\end{choices}
\end{question}
\answer{1}
